%======================================================================%
\section{Wave~2 Thread~I: Emergent Descent Start}
\label{sec:wave2-thread1}
%======================================================================%

Section~\ref{sec:thread-1} records the Thread~1 programme and the evaluator
gap for Class~V physics-map rows: chemistry, atomic spectroscopy, and allied
many-body domains are \emph{mapped} through
$\mathcal{E}_N=(\Pi_{\sigma=-1})^{\otimes N}$ (Definition~\ref{def:emergent-descent})
but lack a named closure functor $\mathcal{F}_{\mathrm{emergent}}$ in
\eqref{eq:physics-map-pipeline}. \textbf{Wave~2} opens Thread~I closure by
certifying the first emergent domain evaluator, proving the descent
factorization for chemistry and atomic rows, and recording the periodic-tower
continuum functor. The results below are \textbf{Tier~B} under ledger entries
\textbf{A5} (flavon overlap truncation) and \textbf{A7} (OS reconstruction on
$\mathcal{H}_N^{\mathrm{em}}$); no new primitive enters beyond
Section~\ref{sec:emergent}.

%----------------------------------------------------------------------%
\subsection{Emergent domain evaluator}
\label{sec:wave2-thread1:evaluator}
%----------------------------------------------------------------------%

Fix the stabilized Albert tower $\mathcal{T}_7$ at $\Mirror^7(\Void)$
(Definition~\ref{def:tower7}), the cellular cosheaf $\mathcal{F}_{16}$, and the
observer collapse $\Pi_{\sigma=-1}$ of Definition~\ref{def:heterotic-collapse}.
Wave~2 assembles the coarse-graining data already present in
Section~\ref{sec:emergent} into a single domain evaluator on $\mathcal{T}_7$.

\begin{definition}[Coarse-graining parameters on $\mathcal{T}_7$]
\label{def:thread1-coarse-params}
For a physics domain $D\in\mathbf{PhysDomains}$, fix:
\begin{enumerate}[label=(\roman*),leftmargin=*]
  \item \textbf{Replica extent $L$.}
    $\mathcal{T}_7^{(L)}$ is the $L$-fold periodic extension of $\mathcal{T}_7$
    along the observer-indexed lattice direction $e_\iota$
    (Definition~\ref{def:replica-limit}).
  \item \textbf{Occupancy $N$.}
    $N\in\mathbb{N}$ is the replicated observer count in
    $\mathcal{H}_N^{\mathrm{em}}=\mathrm{Im}(\mathcal{E}_N)$
    (Definition~\ref{def:emergent-descent}).
  \item \textbf{Inverse temperature $\beta$.}
    $\beta>0$ selects the capacity ensemble $d\nu_\beta$ on
    $\mathcal{H}_{\mathrm{obs}}$ (Definition~\ref{def:capacity-ensemble}).
\end{enumerate}
Write $\mathcal{P}_D=(\beta,N,L)$ for the parameter triple adapted to $D$.
\end{definition}

\begin{definition}[Emergent coarse-graining functor]
\label{def:thread1-coarse}
For $\mathcal{P}_D=(\beta,N,L)$ as in Definition~\ref{def:thread1-coarse-params},
define the \emph{emergent coarse-graining map}
\begin{equation}\label{eq:thread1-coarse}
  \mathrm{coarse}_{\beta,N,L}
  \;:\;
  \mathcal{H}_N^{\mathrm{em}}
  \;\longrightarrow\;
  \mathcal{O}_D,
\end{equation}
where $\mathcal{O}_D$ is the effective observable algebra of $D$:
\begin{itemize}[leftmargin=*]
  \item \textbf{Chemistry / atomic $D$.}
    $\mathcal{O}_D$ is generated by flavon overlap networks $G_Z$
    (Definition~\ref{def:flavon-network}), capacity shell data
    \eqref{eq:shell-ladder}, and the $\beta$-weighted trace on
    $d\mu_{\Gamma_{-1}}$ (Theorem~\ref{thm:thermodynamics}).
  \item \textbf{Condensed-matter $D$.}
    $\mathcal{O}_D$ is generated by Bloch data $\varepsilon_n(\mathbf{k})$
    from \eqref{eq:bloch-bands} on $\mathcal{T}_7^{(L)}$ in the replica limit
    \eqref{eq:replica-limit}.
  \item \textbf{Thermodynamic $D$.}
    $\mathcal{O}_D$ is the Helmholtz sector
    $\{\mathcal{Z}(\beta),F(\beta),S\}$ of Theorem~\ref{thm:thermodynamics}.
\end{itemize}
The map $\mathrm{coarse}_{\beta,N,L}$ sends each tensor factor of
$\mathcal{H}_N^{\mathrm{em}}$ to its cellular locality class on
$\mathcal{T}_7^{(L)}$, projects through $G_Z$ or $\mathcal{D}_{\mathbf{k}}$
as appropriate, and weights expectations by $d\nu_\beta$.
\end{definition}

\begin{definition}[Emergent domain evaluator]\label{def:thread-1-evaluator}
The \emph{emergent domain evaluator} on the Albert tower $\mathcal{T}_7$ is
\begin{equation}\label{eq:thread1-evaluator}
  \mathcal{F}_{\mathrm{emergent}}(D)
  \;:=\;
  \mathrm{coarse}_{\beta,N,L}
  \circ
  \mathcal{E}_N
  \circ
  \Pi_{\sigma=-1},
\end{equation}
with $\mathcal{E}_N=(\Pi_{\sigma=-1})^{\otimes N}$ of
Definition~\ref{def:emergent-descent} and
$\mathrm{coarse}_{\beta,N,L}$ of Definition~\ref{def:thread1-coarse}.
The domain $D$ selects $\mathcal{P}_D=(\beta,N,L)$; for atomic number $Z$,
$N=Z$ and $G_Z$ is the chemistry chart; for spectroscopy,
$\beta$ is the inverse temperature of the capacity ensemble and $L=1$.
The evaluator acts on
$H^1(\mathcal{T}_7,\mathcal{F}_{16})^{\sigma=-1}$ and its $N$-fold tensor
replicas, producing $\mathcal{O}_D$ without postulating a third field species.
\end{definition}

\begin{remark}[Relation to the master pipeline]
Definition~\ref{def:thread-1-evaluator} names the Class~V factor
$\mathcal{F}_D$ in \eqref{eq:physics-map-pipeline}. Composing with the
upstream void--mirror climb gives the Thread~1 boxed pipeline
\eqref{eq:thread-1-pipeline}:
$\mathcal{F}_{\mathrm{emergent}}(D)=\mathrm{coarse}_{\beta,N,L}\circ\mathcal{E}_N
\circ\Pi_{\sigma=-1}\circ\mathrm{Bal}_{\mathcal{C}}\circ\Mirror^{\leq 8}(\Void)$.
Wave~2 certifies the $\mathcal{T}_7$ tail; the Tier~A prefix is unchanged.
\end{remark}

%----------------------------------------------------------------------%
\subsection{Descent factorization for chemistry and atomic domains}
\label{sec:wave2-thread1:descent}
%----------------------------------------------------------------------%

Atomic and chemical physics are not independent of the Standard Model: they are
the coarse-grained image of replicated observer modes under
$\mathcal{L}_{\mathrm{SM}}$ (Definition~\ref{def:sm-lagrangian}).

\begin{definition}[Standard Model sector evaluator]
\label{def:thread1-eval-sm}
Let $\mathfrak{E}$ denote the fundamental contact layer
$\mathcal{U}(\Pi_{\sigma=-1}\circ\mathrm{Bal}_{\mathcal{C}}\circ\Mirror^{\leq 8})
(\Void)$ of Theorem~\ref{thm:universe-summary}. The \emph{Standard Model
sector evaluator} is
\begin{equation}\label{eq:thread1-eval-sm}
  \mathrm{eval}_{\mathrm{SM}}
  \;:\;
  \mathfrak{E}
  \;\longrightarrow\;
  (\mathcal{L}_{\mathrm{SM}},\mathcal{H}_{\mathrm{obs}},\hat H_\iota),
\end{equation}
transporting $\mathfrak{E}$ to the derived SM Lagrangian
\eqref{eq:sm-lagrangian}, the observer Hilbert space, and the OS Hamiltonian
$\hat H_\iota$ of Theorem~\ref{thm:qm-derived}. For a domain $D$, write
$\mathrm{eval}_{\mathrm{SM}}(D)$ for the restriction of
\eqref{eq:thread1-eval-sm} to the sector of $\mathcal{L}_{\mathrm{SM}}$ and
$\mathcal{H}_{\mathrm{obs}}$ relevant to $D$ (electromagnetic atomic
spectroscopy uses $U(1)_{\mathrm{em}}$; chemistry uses
$U(1)_{\mathrm{em}}$--Coulomb plus Pauli structure from replicated fermions).
\end{definition}

\begin{theorem}[Thread~1 descent factorization]\label{thm:thread-1-descent}
Assume Definition~\ref{def:thread-1-evaluator},
Definition~\ref{def:thread1-eval-sm}, ledger entries \textbf{A5} and
\textbf{A7}, and the emergent descent operator
$\mathcal{E}_N$ of Definition~\ref{def:emergent-descent}. For chemistry and
atomic physics domains $D$ in Table~\ref{tab:physics-map-complete}---including
\emph{Chemical bonding}, \emph{Reaction rates / catalysis},
\emph{Atomic spectra}, and \emph{Fine / hyperfine structure}---the emergent
evaluator factorizes as
\begin{equation}\label{eq:thread1-descent}
  \mathcal{F}_{\mathrm{emergent}}(D)
  \;=\;
  \mathrm{eval}_{\mathrm{SM}}(D)
  \circ
  \mathrm{coarse}_{\beta,N,L}
  \circ
  \mathcal{E}_N.
\end{equation}
Equivalently, every Class~V chemistry or atomic observable is the image of
SM+$QED$ kinematics on $\mathcal{H}_N^{\mathrm{em}}$ after cellular
coarse-graining on $\mathcal{T}_7$.
\end{theorem}

\begin{proof}
\emph{Step 1 (many-body seed).}
Definition~\ref{def:emergent-descent} defines
$\mathcal{E}_N=(\Pi_{\sigma=-1})^{\otimes N}$ on
$\mathcal{H}(\mathcal{T}_7,\mathcal{F}_{16})^{\otimes N}$. Proposition~\ref{prop:emergent-functor}
certifies additivity under $N=N_1+N_2$, cellular locality on the three
$\mathbf{16}$-cells, and exclusion of mirror doubles from
$\mathcal{H}_N^{\mathrm{em}}$. For chemistry at atomic number $Z$, set
$N=Z$: each tensor factor is a replicated $\sigma$-anti-invariant $1$-mode on
$\mathcal{T}_7$, and the hub cocycle $\phi_1+\phi_2+\phi_3=0$
(Lemma~\ref{lem:cellular}) couples valence data.

\emph{Step 2 (chemistry coarse-grain).}
Under \textbf{A5}, the flavon overlap truncation
(Definition~\ref{def:flavon-network}) maps
$\mathcal{H}_N^{\mathrm{em}}$ to the bonding graph $G_Z$ with edge weights
$w_{ij}=\sqrt{w_iw_j}\,\mathcal{P}_{ij}$. Theorem~\ref{thm:chemistry-ladder}
identifies shell closures, period rows, group columns, and valence
\eqref{eq:valence} with capacity-saturated layers of replicated
$\Pi_{\sigma=-1}$ modes. Hence
$\mathrm{coarse}_{\beta,N,L}\circ\mathcal{E}_N$ on chemistry domains yields
the overlap-network observable algebra $\mathcal{O}_D$ generated by $G_Z$ and
the shell ladder \eqref{eq:shell-ladder}.

\emph{Step 3 (atomic coarse-grain).}
For atomic spectroscopy, $L=1$ and $\beta$ selects the capacity ensemble
\eqref{eq:capacity-ensemble}. Theorem~\ref{thm:thermodynamics}, item~(i),
identifies $\mathcal{Z}(\beta)$ with the $\Pi_{\sigma=-1}$-projected
functional integral on $\mathcal{M}_{-1}$. Coarse-graining therefore
produces $\beta$-weighted expectation values of electromagnetic observables on
$\mathcal{H}_N^{\mathrm{em}}$ without importing a non-SM Hamiltonian.

\emph{Step 4 (SM factor).}
Theorem~\ref{thm:qm-derived}, under \textbf{A7}, reconstructs unitary quantum
mechanics on $\mathcal{H}_{\mathrm{obs}}$ with Hamiltonian $\hat H_\iota$.
Definition~\ref{def:sm-lagrangian} and Theorem~\ref{thm:sm-lagrangian-derived}
supply $\mathcal{L}_{\mathrm{SM}}$ including $U(1)_{\mathrm{em}}$ and three
generations of fermions. Atomic spectra and fine/hyperfine structure are
$U(1)_{\mathrm{em}}$--QED matrix elements of $\hat H_\iota$ on
$\mathcal{H}_N^{\mathrm{em}}$; chemical bonding is the low-energy Coulomb plus
Pauli avatar of the same replicated fermion data. The map
$\mathrm{eval}_{\mathrm{SM}}(D)$ extracts these matrix elements from
$\mathfrak{E}$, yielding \eqref{eq:thread1-descent}.

\emph{Step 5 (factorization).}
Steps~1--3 construct
$\mathrm{coarse}_{\beta,N,L}\circ\mathcal{E}_N$; Step~4 shows
$\mathrm{eval}_{\mathrm{SM}}(D)$ is the unique SM+$QED$ lift of the coarse
data. No intermediate primitive enters; the composition is functorial under
the additivity of Proposition~\ref{prop:emergent-functor}, proving
\eqref{eq:thread1-descent}.
\end{proof}

\begin{remark}[Tier assignment]
Theorem~\ref{thm:thread-1-descent} is \textbf{Tier~B}: Step~2 cites
\textbf{A5}; Steps~3--4 cite \textbf{A7} and the SM Lagrangian certificate
(Theorem~\ref{thm:sm-lagrangian-derived}). Shell structure and valence in
Step~2 are cellular consequences of observer collapse (Tier~A); quantitative
ionization energies and sub-percent spectroscopy remain \textbf{Tier~C}
contacts per Remark~\ref{rmk:emergent-honest}.
\end{remark}

%----------------------------------------------------------------------%
\subsection{Periodic-tower continuum functor}
\label{sec:wave2-thread1:continuum}
%----------------------------------------------------------------------%

Condensed-matter and hydrodynamic Class~V rows require a named continuum limit
of the cellular cosheaf on $\mathcal{T}_7^{(L)}$.

\begin{definition}[Continuum limit functor]\label{def:thread-1-continuum}
Let $\{\mathcal{T}_7^{(L)}\}_{L\ge 1}$ be the nested family of periodic towers
obtained by translating the three-cell star of Definition~\ref{def:tower7}
along the observer-indexed lattice direction $e_\iota$, with cellular cosheaf
$\mathcal{F}_{16}$ restricted to each level. The \emph{continuum limit
functor} is
\begin{equation}\label{eq:thread1-continuum}
  \mathcal{F}_{\mathrm{cont}}
  \;:=\;
  \varinjlim_{L\to\infty}
  \Bigl(
    H^1\!\bigl(\mathcal{T}_7^{(L)},\mathcal{F}_{16}\bigr)^{\sigma=-1},
    \;\mathcal{D}_{\mathbf{k}}^{(L)}
  \Bigr),
\end{equation}
where $\mathcal{D}_{\mathbf{k}}^{(L)}$ is the capacity-weighted Bloch
connection \eqref{eq:bloch-bands} on $\mathcal{T}_7^{(L)}$ and the colimit is
taken in the replica thermodynamic limit
\begin{equation}\label{eq:thread1-replica}
  L\to\infty,
  \qquad
  N=\rho\,L^3,
  \qquad
  \rho>0\ \text{fixed}
\end{equation}
of Definition~\ref{def:replica-limit}. The output is the effective continuum
action
\begin{equation}\label{eq:thread1-cont-action}
  S_{\mathrm{eff}}^{\mathrm{cont}}
  \;=\;
  \int d^4x\,\sqrt{-g}\,
  \Bigl[
    \sum_n\int_{\mathrm{BZ}}\frac{d^3k}{(2\pi)^3}\,
    f\!\bigl(\varepsilon_n(\mathbf{k})-\mu\bigr)
    \;+\;
    \mathcal{L}_{\mathrm{hydro}}^{\mathrm{eff}}
  \Bigr],
\end{equation}
with band energies $\varepsilon_n(\mathbf{k})$ from \eqref{eq:bloch-bands},
Fermi--Dirac occupancy $f$, and $\mathcal{L}_{\mathrm{hydro}}^{\mathrm{eff}}$
the hydrodynamic envelope of $\mathcal{H}_N^{\mathrm{em}}$ under
$\mathrm{coarse}_{\beta,N,L}$ at fixed density $\rho$.
\end{definition}

\begin{proposition}[Continuum functoriality]
\label{prop:thread1-continuum-functor}
The assignment $L\mapsto\mathcal{T}_7^{(L)}$ and the colimit
\eqref{eq:thread1-continuum} satisfy:
\begin{enumerate}[label=(\roman*),leftmargin=*]
  \item \textbf{Dimension stability.}
    $\dim H^1(\mathcal{T}_7^{(L)},\mathcal{F}_{16})^{\sigma=-1}=3$ for all
    $L$, by replication of Lemma~\ref{lem:cellular}.
  \item \textbf{Band convergence.}
    $\varepsilon_n^{(L)}(\mathbf{k})\to\varepsilon_n(\mathbf{k})$ as
    $L\to\infty$ in the norm of Theorem~\ref{thm:condensed-matter},
    item~(ii).
  \item \textbf{Density-fixed extensivity.}
    With $N=\rho L^3$, the free energy per unit volume
    $F/N$ converges to the band-theory potential
    $\Omega(\mu,T)$ of Theorem~\ref{thm:thermodynamics}, item~(iv).
\end{enumerate}
\end{proposition}

\begin{proof}
\emph{(i)} Each periodic cell is a translate of the three-cell star; the
$\sigma=-1$ cohomology dimension is a cellular count unchanged under
translation (Lemma~\ref{lem:cellular}).

\emph{(ii)} Theorem~\ref{thm:condensed-matter} identifies Bloch bands as
eigenvalues of $\mathcal{D}_{\mathbf{k}}$ propagated along
$\mathcal{T}_7^{(L)}$; increasing $L$ refines the Brillouin-zone sampling
without altering the local connection on $\mathcal{F}_{16}$, so band energies
converge.

\emph{(iii)} Theorem~\ref{thm:thermodynamics}, Step~4, identifies the
$N\to\infty$ limit of $\mathcal{Z}(\beta)$ with the band determinant of
$\mathcal{D}_{\mathbf{k}}$; the density-fixed scaling $N=\rho L^3$ in
\eqref{eq:thread1-replica} yields extensivity of $F/N$.
\end{proof}

\begin{remark}[Tier assignment: continuum]
Definition~\ref{def:thread-1-continuum} is \textbf{Tier~B/C}: the replica
limit and band convergence are Tier~B under \textbf{A7} and the emergent
lattice identification of Definition~\ref{def:replica-limit}; quantitative band
gaps and hydrodynamic transport coefficients are \textbf{Tier~C} contacts.
Definition~\ref{def:thread-1-continuum} supplies the continuum functor named in
Table~\ref{tab:thread-1-questions} row EM-FLUID; full Navier--Stokes closure
awaits a later wave.
\end{remark}

%----------------------------------------------------------------------%
\subsection{Wave~2 partial closure certificate}
\label{sec:wave2-thread1:closure}
%----------------------------------------------------------------------%

\begin{corollary}[Wave~2 Thread~I partial closure]
\label{cor:thread-1-wave2-partial}
Relative to Section~\ref{sec:thread-1}, Table~\ref{tab:thread-1-questions},
and Class~V of Section~\ref{sec:physics-map:grouping}:
\begin{enumerate}[label=(\roman*),leftmargin=*]
  \item \textbf{Evaluator certified.}
    Definition~\ref{def:thread-1-evaluator} names the first closed Class~V
    functor $\mathcal{F}_{\mathrm{emergent}}$ in
    \eqref{eq:physics-map-pipeline}, assembling
    $\mathrm{coarse}_{\beta,N,L}\circ\mathcal{E}_N\circ\Pi_{\sigma=-1}$ on
    $\mathcal{T}_7$.
  \item \textbf{Descent factorization.}
    Theorem~\ref{thm:thread-1-descent} proves
    $\mathcal{F}_{\mathrm{emergent}}(D)=\mathrm{eval}_{\mathrm{SM}}(D)\circ
    \mathrm{coarse}_{\beta,N,L}\circ\mathcal{E}_N$ for chemistry and atomic
    domains under \textbf{A5} and \textbf{A7}.
  \item \textbf{Continuum functor.}
    Definition~\ref{def:thread-1-continuum} records the $L\to\infty$ periodic
    tower colimit $\mathcal{F}_{\mathrm{cont}}$ on
    $\mathcal{T}_7^{(L)}$, with band and extensivity certificates of
    Proposition~\ref{prop:thread1-continuum-functor}.
  \item \textbf{Physics-map upgrade.}
    The following rows of Table~\ref{tab:physics-map-complete} now carry the
    named evaluator $\mathcal{F}_D=\mathcal{F}_{\mathrm{emergent}}$ rather than
    prose-only $\mathrm{eval}_{\mathrm{emergent}}$:
    \begin{center}
    \small
    \renewcommand{\arraystretch}{1.08}
    \begin{tabular}{@{}lll@{}}
    \toprule
    \textbf{Domain} & \textbf{Subfield} & \textbf{Functor} \\
    \midrule
    Molecular & Chemical bonding &
      $\mathcal{F}_{\mathrm{emergent}}$ via $G_Z$
      (Theorem~\ref{thm:chemistry-ladder}) \\
    Chemistry & Reaction rates / catalysis &
      $\mathcal{F}_{\mathrm{emergent}}$ via $G_Z$, $d\nu_\beta$ \\
    Atomic & Atomic spectra &
      $\mathcal{F}_{\mathrm{emergent}}$ via
      Theorem~\ref{thm:thread-1-descent} \\
    Atomic & Fine / hyperfine structure &
      $\mathcal{F}_{\mathrm{emergent}}$ via $\mathrm{eval}_{\mathrm{SM}}(D)$,
      QED on $U(1)_{\mathrm{em}}$ \\
    \bottomrule
    \end{tabular}
    \end{center}
\end{enumerate}
Residual Thread~1 items---band-gap contacts (\texttt{thm:thread-1-bands}),
critical exponents (\texttt{thm:thread-1-critical}), strong-correlation phases
(\texttt{prop:thread-1-correlation}), and full Class~V closure
(\texttt{cor:thread-1-closure})---remain open for subsequent waves.
\end{corollary}

\begin{proof}
Item~(i) is Definition~\ref{def:thread-1-evaluator}. Item~(ii) is
Theorem~\ref{thm:thread-1-descent}. Item~(iii) is
Definition~\ref{def:thread-1-continuum} and
Proposition~\ref{prop:thread1-continuum-functor}. Item~(iv) follows from
\eqref{eq:thread1-descent} and the row citations in
Table~\ref{tab:physics-map-complete}: chemistry and atomic subfields factor
through $\mathcal{F}_{\mathrm{emergent}}$ with explicit $\mathcal{T}_7$
data, upgrading the Class~V grouping of
Section~\ref{sec:physics-map:grouping} from Tier~E contact to Tier~B
evaluator status for these four rows.
\end{proof}

\begin{remark}[Thread~1 question ledger upgrade]
\label{rmk:thread1-wave2-table}
After Wave~2, Table~\ref{tab:thread-1-questions} updates as follows:
\begin{center}
\small
\begin{tabular}{@{}clc@{}}
\toprule
\textbf{ID} & \textbf{Resolution status} & \textbf{Tier} \\
\midrule
EM-EVAL &
  \textbf{Closed} --- Definition~\ref{def:thread-1-evaluator};
  $\mathcal{F}_{\mathrm{emergent}}$ certified &
  B \\
EM-CHEM &
  \textbf{Partial} --- Theorem~\ref{thm:thread-1-descent} + shell ladder;
  spectroscopic refinement open &
  B/C \\
EM-ATOM &
  \textbf{Partial} --- Theorem~\ref{thm:thread-1-descent};
  \texttt{cor:thread-1-qed} contact layer open &
  B/C \\
EM-FLUID &
  \textbf{Started} --- Definition~\ref{def:thread-1-continuum};
  hydrodynamic closure open &
  B/C \\
\bottomrule
\end{tabular}
\end{center}
Entries EM-BAND, EM-SC, EM-THERMO, EM-ASTRO, and EM-BIO remain at their
Section~\ref{sec:thread-1} programme status pending later waves.
\end{remark}